% CP-VI-0064
% target_chapter: VI/12 — Presheaves
% solution_provenance: FRESH_CANONICAL_AUTHORED_SOLUTION

\subsection*{CP-VI-0064: A presheaf that is not a sheaf}

On a topological space $X$ with at least two disjoint nonempty open sets, define $F(U)=\mathbb Z$ for every nonempty open $U$, $F(\varnothing)=0$, and use identity restriction maps between nonempty opens. Show that $F$ is a presheaf but generally not a sheaf.

\medskip
\noindent\textbf{Solution.}

Identity maps satisfy the presheaf identities. Let $U=U_1\sqcup U_2$ with $U_1,U_2$ disjoint and nonempty. Choose sections $s_1=0\in F(U_1)$ and $s_2=1\in F(U_2)$. Their restrictions agree on the empty overlap. A sheaf would provide $s\in F(U)=\mathbb Z$ restricting simultaneously to $0$ and $1$, but both restrictions are the identity, impossible. Hence the gluing axiom fails.
